% ===================================================================
%  REGISTRO DE CAMBIOS DE CONSISTENCIA CON LA CONSIGNA
%  Proyecto Global Integrador AyME 2026 — Calderón / Castel
%  Generado: 12/06/2026
%  Compilar con: pdflatex registro_cambios_consigna.tex
% ===================================================================
\documentclass[11pt,a4paper]{article}
\usepackage[utf8]{inputenc}
\usepackage[T1]{fontenc}
\usepackage{lmodern}
\usepackage[spanish]{babel}
\usepackage{amsmath,amssymb,bm}
\usepackage[margin=2.2cm]{geometry}
\usepackage{xcolor}
\usepackage{longtable}
\usepackage{booktabs}
\usepackage{enumitem}
\definecolor{okgreen}{RGB}{0,128,0}
\definecolor{warnorange}{RGB}{204,102,0}
\definecolor{errred}{RGB}{180,0,0}
\newcommand{\hecho}{\textcolor{okgreen}{\textbf{HECHO}}}
\newcommand{\pendiente}{\textcolor{errred}{\textbf{PENDIENTE}}}
\newcommand{\evaluar}{\textcolor{warnorange}{\textbf{A EVALUAR}}}

\title{Registro de cambios de consistencia con la consigna\\
\large Informe Técnico (\texttt{main.tex}) y Presentación (\texttt{main\_presentacion.tex})}
\author{Proyecto Global Integrador --- AyME 2026}
\date{12 de junio de 2026}

\begin{document}
\maketitle

\section*{Resumen}
Este documento registra, uno por uno, los cambios de nomenclatura y contenido aplicados al informe y a la presentación para alinearlos con la \emph{Guía de Trabajo y Especificaciones} del Ing.\ G.\,L.\ Julián (Rev.\,0), tras la auditoría exhaustiva de 35 discrepancias verificadas. Se listan también los ítems \textbf{pendientes} (trabajo de simulación aún no realizado) y los ítems \textbf{a evaluar manualmente}.

Los IDs corresponden a la numeración de la auditoría. El punto de partida de la comparación es la versión del informe subida por el compañero de equipo (respaldada como \texttt{aux\_main\_anterior.tex}); el punto final es la versión actual de ambos documentos. La Sección~3 registra además los cambios de método, redacción y presentación posteriores al lote principal. Ambos documentos compilan sin errores tras los cambios (informe: 117 páginas; presentación: 186 páginas).

% ===================================================================
\section{Cambios realizados --- correcciones previas (sesión anterior)}

\begin{enumerate}[label=\textbf{D\arabic*.}]
  \item \hecho\ $n_{l,nom}$: corregido de 55 a \textbf{60 rpm} en la tabla de parámetros (valor de la consigna).
  \item \hecho\ Superíndice rotórico: $i_{qs}\to i_{qs}^r$, $i_{ds}\to i_{ds}^r$ en los modelos LTI nominal y aumentado, FT (incl.\ $I_{qs}^r(s)$), observabilidad, simulación y \emph{abstract} (70 reemplazos).
  \item \hecho\ Flujo de imanes unificado a $\lambda'^r_m$ (49 reemplazos; eliminadas las variantes $\lambda'_m$, $\lambda_m'$, $\lambda_m$ y $\lambda_r$).
  \item \hecho\ $R_{sREF}$ unificado (eliminadas $R_{s,REF}$ y $R_{s,ref}$); además $R_{ts,amb}\to R_{ts\text{-}amb}$.
  \item \hecho\ $\alpha_{Cu}$ unificado (eliminada $\alpha_{cu}$).
  \item \hecho\ Símbolo de grado agregado: $T_s^\circ(t)$, $T_{amb}^\circ(t)$ en LPV, LTI, tabla y conclusiones (30 reemplazos).
\end{enumerate}

\textbf{IDs 9 y 10} (corregidos en esta sesión, antes del lote principal):
\begin{itemize}
  \item \hecho\ \textbf{ID 9}: $T_{REF}^\circ$ y $T_{s,ref}$ unificados a $T_{sREF}^\circ$ (Ec.\ de $R_s$, línea 383; migración ante $R_s$; presentación).
  \item \hecho\ \textbf{ID 10 (error de física)}: la fórmula de la sección «Migración ante variación de $R_s$» mezclaba $R_{s0}=1{,}058\ \Omega$ (valor a $29{,}5\,^\circ$C) con la referencia de $20\,^\circ$C. Quedó la ley física oficial $R_s(T_s^\circ)=R_{sREF}\,[1+\alpha_{Cu}(T_s^\circ - T_{sREF}^\circ)]$ (Ec.\ 3.9 de la guía), aclarando que el \emph{barrido} se hace alrededor del valor de operación $R_{s0}$. Nota: $R_{s0}$ se mantiene en el resto del documento por ser un símbolo deliberado del modelo LTI (definido en las secciones 5.1.2.c y 5.1.6), con 42 usos consistentes.
\end{itemize}

% ===================================================================
\section{Cambios realizados --- lote principal (esta sesión)}

\begin{enumerate}[label=\textbf{ID \arabic*.}, leftmargin=1.6cm]
  \item \hecho\ Declarada la equivalencia $q(t)\equiv\theta_l(t)$ como \textbf{variable controlada} (no medida directamente), con $q=\theta_l=\frac{1}{r}\theta_m$, en la Sección 3.1 del informe.
  \item \hecho\ Declarada la equivalencia $\tau(t)\equiv T_q(t)$ del torque impulsor (Sección 3.1).
  \item \hecho\ Declarada la referencia angular de la consigna: $\theta_l$ medida desde el eje vertical hacia abajo (equilibrio estable), positiva antihoraria; $g$ como perturbación externa vertical constante (Sección 3.1).
  \item \hecho\ Agregada la nota del encoder según la guía: encoder incremental con \emph{homing} y decodificación idealizados; $\theta_m(t)$ absoluta «rectificada»; ajuste $\theta_m\equiv 0$ (\emph{Home}) para $\theta_l\equiv 0$ (sección de observabilidad).
  \item \hecho\ Definido el \emph{setpoint} de la guía $q_1^*(t)\equiv\frac{1}{r}\theta_m^*(t)$ como entrada de referencia del lazo externo (Sección 5.2.2).
  \item \hecho\ Declarado el perfil trapezoidal de la guía en 5.2.4.a: $q_1^*(t)=0 \xrightarrow{\Delta t_{ramp}=5\,\text{s}} 2\pi\ \text{rad} \xrightarrow{\Delta t_{ramp}=5\,\text{s}} 0$.
  \item \hecho\ $T_{Ld}\to T_{ld}$ (subíndice en minúscula, 3+3 instancias en informe y presentación).
  \item \hecho\ $P_{Cu}\to P_{s,perd}$ (notación oficial de la Ec.\ 3.10 de la guía; 3+4 instancias).
  \setcounter{enumi}{10}
  \item \hecho\ \textbf{Puntos de operación unificados al estilo de la consigna} ($I^r_{ds,o}$): todas las variantes ($i_{qsO}^r$, $I_{ds0}^r$, $\omega_{mO}$, $\theta_{mO}$, $T_{sO}$, $T_{ambO}$, $R_{sO}$, $v_{qsO}^r$, $\Gamma_{th\,O}$, \dots) pasaron a subíndice «$,o$» minúscula: $i^r_{qs,o}$, $I^r_{ds,o}$, $\omega_{m,o}$, $\theta_{m,o}$, $T_{s,o}$, $T_{amb,o}$, $R_{s,o}$, $v^r_{qs,o}$, $\Gamma_{th,o}$ ($\sim$75+50 reemplazos en informe y presentación).
  \item \hecho\ Secuencia cero \textbf{sin} superíndice $r$ (la componente $0$ no pertenece al marco rotórico, Ec.\ 3.8): $i_{0s}^r\to i_{0s}$, $v_{0s}^r\to v_{0s}$ en modelo NL, LPV y demostración de Park (22+10 instancias).
  \item \hecho\ \textbf{Eliminado el símbolo $V_m$} (no existe en la consigna): las tensiones de fase quedaron con la forma oficial $v_{as}(t)\approx\sqrt{2}\,\frac{V_{sl}(t)}{\sqrt{3}}\cos(\theta_{ev}(t))$ (Ecs.\ 4.1--4.3), y la derivación de Clarke/atan2 usa la amplitud $\sqrt{2}\,V_{sl}/\sqrt{3}$ explícita.
  \item \hecho\ Declaradas las \textbf{variables manipulables del inversor} según la guía: $V_{sl}(t)\in[0;48]$ Vca rms y $f_e(t)\in[-330;+330]$ Hz (su signo define la secuencia $abc/acb$ y el sentido de giro), con $\omega_e(t)\equiv 2\pi f_e(t)\equiv \frac{d\theta_{ev}}{dt}$.
  \item \hecho\ La saturación del inversor ahora muestra el valor de la consigna: $|v_{abc}|\le\frac{\sqrt{2}\cdot V_{sl,max}}{\sqrt{3}}$ con $V_{sl,max}=48$ Vca rms $\Rightarrow 39{,}19$ V.
  \item \hecho\ \textbf{Tabla de parámetros completada} con los límites de la consigna: $n_{m,nom}=6600$ rpm ($\omega_{m,nom}=691{,}15$ rad/s), $V_{sl,nom}=30$ Vca rms, $V_{sl,max}=48$ Vca rms, $I_{s,nom}=0{,}4$ Aca rms, $I_{s,max}=2{,}0$ Aca rms, $T^\circ_{s,max}=115\,^\circ$C, rango $T^\circ_{amb}\in[-15;+40]\,^\circ$C, $\omega_{l,nom}=6{,}28$ rad/s, y nota de tolerancia $\pm 1\%$. Además, en 5.1.6 se agregó la verificación de que la velocidad de vacío ($\approx 405$ rad/s) es inferior al límite $\omega_{m,nom}=691{,}15$ rad/s.
  \item \hecho\ $b_l = 0{,}1 \pm 0{,}03$ Nms/rad (incertidumbre de la consigna) y nueva fila $T_{ld}=0\pm 5{,}0$ Nm (escalón) en la tabla.
  \item \hecho\ Agregadas las fórmulas de composición de la guía con sus rangos: $J_l=(m\,l_{cm}^2+J_{cm})+m_l\,l_l^2 = 0{,}0833+[0\ldots 0{,}375]$ kg·m$^2$ y $k_l=m\,l_{cm}+m_l\,l_l = 0{,}25+[0\ldots 0{,}75]$ kg·m (Sección 3.1 + tabla).
  \setcounter{enumi}{19}
  \item \hecho\ $T_{s,\max}\to T^\circ_{s,max}$ (símbolo de grado, 2+3 instancias).
  \item \hecho\ $T_{em}\to T_m$ en toda la sección 5.1.6 (único símbolo de la Ec.\ 3.5 de la guía; 7+3 instancias).
  \setcounter{enumi}{22}
  \item \hecho\ Ganancias del observador unificadas: $K_{e\theta_m}$, $K_{e\omega_m}$, $K_{ei_m}$ en 5.2.3 y 5.2.5.b (eliminadas las variantes $K_{e_{\theta_m}}$, $K_{e\theta}$, $K_{e\omega}$, $K_{ei}$), más nota en 5.2.7.b con la correspondencia a las variables del código (\texttt{K\_e\_tita\_m\_int}, \texttt{K\_e\_w\_m\_int}, \texttt{K\_e\_i\_m\_int}).
  \item \hecho\ Polos con el nombre de la consigna: lazos de corriente siempre $p_i=-5000$ rad/s (eliminadas $p_{1,2,3}$ y $p_{a,b,c}$); observador siempre $p_{obs\,1,2}=-3200$ rad/s y $p_{obs\,1,2,3}$ para el aumentado (los $p_{1,2}$ de los \emph{sensores} se mantienen, así los nombra la consigna).
  \setcounter{enumi}{25}
  \item \hecho\ Nodo neutro renombrado de $n$ a $O$ («centro de estrella $O$», como la consigna), incl.\ figura circuital y $v_{an}\to v_{aO}$.
  \item \hecho\ Instantes de conmutación renombrados $t_k \to t_{step\,k}$ (notación de la consigna, valores idénticos).
  \item \hecho\ Consignas de tensión de 5.1.6 con superíndice rotórico: $v_{qs}^*\to v_{qs}^{r*}$, $v_{ds}^*\to v_{ds}^{r*}$, $V_{ds}^*\to V_{ds}^{r*}$, $V_{qs,nom}\to V_{qs,nom}^{r}$ (24+15 y 10+8 instancias).
  \setcounter{enumi}{32}
  \item \hecho\ $\Gamma_{th}$ unificada a $\Gamma_{th,o}$ (matriz, texto y definición).
  \item \hecho\ $R_{ts\text{-}amb}$ con guion tipográfico en todas las instancias (el guion matemático se renderizaba como resta).
  \item \hecho\ Unidades de la guía en la tabla: $\lambda'^r_m$ en \textbf{Wb-turn}; $C_{ts}$ en \textbf{W/($^\circ$C/s)}.
\end{enumerate}

% ===================================================================
\section{Cambios realizados --- método, redacción y presentación (posteriores al lote principal)}

\begin{enumerate}[label=\textbf{P\arabic*.}, leftmargin=1.3cm]
  \item \hecho\ \textbf{Reescritura completa del diseño del controlador externo (sintonía serie bien fundamentada).} La versión anterior definía a priori tres frecuencias características ($\omega_{vel}:=b_a/J_{eq}$, $\omega_{pos}:=K_{sa}/b_a$, $\omega_{int}:=K_{sia}/K_{sa}$) y un parámetro de separación $n$ ($\omega_{vel}=n\,\omega_{pos}$, $\omega_{int}=\omega_{pos}/n$) del que luego se deducía $n=2\zeta+1$, sin justificar el origen de las relaciones. Se reemplazó por el desarrollo canónico de la cátedra:
  \begin{enumerate}[label=(\roman*)]
      \item dinámica de lazo cerrado del PID sobre la planta equivalente (doble integrador $J_{eq}\ddot\theta_m = T_m^{*\prime} - T_{ld}/r$, con lazos de corriente $6{,}25\times$ más rápidos y fricción/gravedad precompensadas), con las dos funciones de transferencia y el rechazo exacto de perturbación constante ($G_{T_{ld}}(0)=0$);
      \item polinomio \emph{estructural} $P_e(s) = s^3+\frac{b_a}{J_{eq}}s^2+\frac{K_{sa}}{J_{eq}}s+\frac{K_{sia}}{J_{eq}}$;
      \item polinomio \emph{requerido} $P_r(s)=(s+\omega_n)(s^2+2\zeta\omega_n s+\omega_n^2)$ expandido, donde el factor $(2\zeta+1)$ surge como identidad algebraica;
      \item igualación coeficiente a coeficiente y despeje directo: $b_a=J_{eq}\,\omega_n(2\zeta+1)$, $K_{sa}=J_{eq}\,\omega_n^2(2\zeta+1)$, $K_{sia}=J_{eq}\,\omega_n^3$.
  \end{enumerate}
  \textbf{Los valores numéricos de las ganancias no cambian} ($0{.}0396$; $31{.}656$; $10\,129{.}778$). La derivación fue verificada en forma independiente (álgebra paso a paso + cálculo numérico): sin errores; los polos resultantes son $-800$ y $-600\pm j\,529{,}15$, los tres de módulo $\omega_n$.
  \item \hecho\ Se agregó la \textbf{figura de la configuración de polos} asignada por $P_r(s)$ (estilo de la cátedra): polo real en $-\omega_n$ y par conjugado de parámetros $\zeta$, $\omega_n$ sobre la circunferencia punteada de radio $\omega_n$, en forma \textbf{genérica} (sin valores numéricos de los polos), como ilustración del método.
  \item \hecho\ Se eliminó por completo el desarrollo del parámetro $n$ (frecuencias postuladas, jerarquía anidada y nota de vínculo). El método se denomina \emph{sintonía serie} desde la introducción de la sección, con redacción sobria y sin sobre-explicar la consigna.
  \item \hecho\ \textbf{Presentación espejada}: slides nuevas de dinámica de lazo cerrado, polinomios $P_e$/$P_r$, figura de polos genérica, igualación/despeje y valores numéricos; eliminada la slide del vínculo $n=2\zeta+1$; el título de la slide de polinomios lleva el nombre del método («Método de sintonía serie») y la slide de objetivo indica «Método: sintonía serie».
  \item \hecho\ \textbf{Eliminadas todas las referencias a los puntos de la guía en el texto del informe} (10 reescrituras): los títulos «(Punto 5.2.5.b)» y «(Punto 5.2.5.c)»; «para el punto 5.2.3»; «el comportamiento exigido en el punto 5.2.5.b»; «correspondiente al punto 5.2.3 \dots\ para el punto 5.2.5.b»; «solución final para el punto 5.2.5.b»; «diseñado en la sección~5.2.5.b»; «en la sección 5.2»; «de/en 5.1.2.c.VI» (título y texto del párrafo de cambio de dinámica); «, en 5.2.2,»; y «(Ec.~3.9 de la guía)». Todas sustituidas por redacción natural («más adelante», «previamente», «en la verificación de desempeño», etc.). Verificación final por patrones: 0 ocurrencias remanentes en el informe. En la presentación las únicas menciones están en comentarios del código fuente (no se renderizan) y se conservaron como navegación.
  \item \hecho\ \textbf{Limpieza de la presentación}: se eliminaron los 109 bloques de texto gris «ayuda-memoria» que el template nunca renderizó (Beamer los absorbía como subtítulo de \emph{frame} y la plantilla personalizada no muestra subtítulos), y la macro \texttt{\textbackslash figslide} pasó de 4 a 3 argumentos (61 llamadas adaptadas, se quitó la descripción muerta). Verificado con \texttt{pdftotext}: el texto del PDF es idéntico línea por línea antes y después.
\end{enumerate}

% ===================================================================
\section{Ítems PENDIENTES (trabajo de simulación no realizado)}

\begin{enumerate}[label=\textbf{ID \arabic*.}, leftmargin=1.6cm]
  \setcounter{enumi}{18}
  \item \evaluar\ \textbf{Estado inicial consistente y peor caso térmico.} La consigna (4.c) da como ejemplo el estado inicial $\theta_l(t_0)=0$, $\omega_m(t_0)=0$, $\bm{i}^r_{qd0s}(t_0)=\bm{0}$, $T^\circ_s(t_0)=T^\circ_{amb}(t_0)=T^\circ_{amb,max}=40\,^\circ$C. El informe usa $20\,^\circ$C (validación) y $25\,^\circ$C (resto) sin declarar el estado inicial de la guía ni justificar la elección. En particular, la verificación térmica 5.2.5.c nunca se evalúa en el peor caso $T_{amb}=40\,^\circ$C (donde el límite de $115\,^\circ$C se alcanzaría antes de los $225{,}9$ s reportados). \emph{Decisión del equipo: evaluar más adelante.}
  \setcounter{enumi}{28}
  \item \pendiente\ \textbf{Rechazo a perturbaciones con variación máxima de parámetros (5.2.4.b).} La consigna pide simular el rechazo a escalones de $T_{ld}$ con valores nominales \textbf{y} con la variación máxima de los parámetros de carga ($m_l=1{,}5$ kg; $b_l=0{,}1\pm0{,}03$), manteniendo el control invariante. Solo está hecho el caso nominal. Falta la corrida WORST-CASE en Simulink y su comparación. (Las propias Conclusiones lo admiten como pendiente.)
  \item \pendiente\ \textbf{Comparación de estrategias del modulador (Nota de 5.2.1).} La consigna pide \emph{simular} con el Modulador de Torque en DT y comparar contra el comportamiento de 5.1.6, contrastando la estrategia base $i^{r*}_{ds}(t)\equiv 0$ vs.\ la extendida $i^{r*}_{ds}(t)\neq 0$ (debilitamiento/reforzamiento). Hoy solo hay una discusión conceptual de ambos casos; falta la simulación comparativa.
  \item \pendiente\ \textbf{Verificación de límites del reductor (5.2.5.a).} La sección de verificación solo analiza corriente (motor) y tensión (inversor). Falta verificar numéricamente la velocidad de salida $\omega_l$ contra $n_{l,nom}=60$ rpm y el torque $T_q$ contra $T_{q,nom}=17$ / $T_{q,max}=45$ N·m.
  \item \pendiente\ \textbf{Gráficas faltantes de 5.1.6.} La consigna pide graficar tensiones y corrientes \emph{en ambas} coordenadas ($qd0 \Leftrightarrow abcs$) y superponer las curvas características cuasi-estacionarias sobre la curva paramétrica $T_m$--$\omega_m$. Faltan: la gráfica de \textbf{tensiones en $abc$} ($v_a$, $v_b$, $v_c$) y la comparación con las \textbf{curvas características cuasi-estacionarias}.
\end{enumerate}

% ===================================================================
\section{Ítems A EVALUAR MANUALMENTE (detallados, sin cambios aplicados)}

\subsection*{ID 22 --- Cuatro nombres para la consigna de torque del PID \evaluar}
La misma señal (la consigna de torque que entrega el lazo externo) aparece con \textbf{cuatro} nomenclaturas distintas según la sección:
\begin{itemize}
  \item $T_m^{\ast\prime}(t)$ --- en el \emph{precomputed torque} del modulador (5.2.1) y en la ley del PID (5.2.2): $T_m^{*\prime}(s)=b_a e_\omega + K_{sa}e_{\theta_m} + K_{sia}e_{\theta_m}/s$.
  \item $T_m'(t)$ --- en el observador continuo (5.2.3 y 5.2.5.b): $u(t)=T_m'(t)$, $\dot{\hat\omega}_m = \frac{1}{J_{eq}}T_m' + \dots$
  \item $T_m^{**}$ --- en la introducción de la codificación (5.2.7): «el controlador PID de posición calcula una referencia de torque mecánico $T_m^{**}$».
  \item $T_e^{*}$ --- en el observador codificado (5.2.7.b): «recibe \dots el torque electromagnético de referencia $T_e^{*}$» (variable \texttt{Te\_ref}).
\end{itemize}
\textbf{Por qué no se cambió:} la unificación automática es riesgosa: (a) hay una diferencia \emph{conceptual} sutil entre la salida pura del PID ($T_m^{*\prime}$) y el torque electromagnético de referencia que entra al observador ($T_e^*$, que en rigor es $T_m^*$ después del precómputo); unificar mal podría introducir un error técnico; (b) cambiar $T_e^*/$\texttt{Te\_ref} obligaría a tocar el código MATLAB y sus capturas de pantalla (\texttt{codificado\_observador.png}), que quedarían desincronizadas. \textbf{Sugerencia para evaluar:} decidir primero qué señal entra físicamente al observador (¿$T_m^{*\prime}$ o $T_m^{*}$?) y luego unificar texto + código + figuras en una sola pasada coherente.

\subsection*{ID 25 --- Ángulo de torque del rotor \evaluar}
Tres diferencias con la consigna (sección 3.6 de la guía) en el bloque del ángulo $\delta(t)$ de 5.1.6:
\begin{enumerate}
  \item \textbf{Nombre:} la guía lo llama «ángulo de \emph{torque} del rotor»; el informe lo llama «Ángulo de \emph{Carga}». (Son sinónimos técnicos --- \emph{load angle} --- pero la guía pide consistencia.)
  \item \textbf{Errata en la ecuación:} el informe escribe $\delta(t) = \int_0^t[\omega_e-\omega_r]d\xi + [\theta_e(0)-\theta_r(0)]$; el término inicial debería ser $\theta_{ev}(0)$, no $\theta_e(0)$ (el subíndice $e$ a secas no está definido).
  \item \textbf{Interpretación faltante:} la guía pide enunciar que el rotor queda «atrasado» en modos de motorización y «adelantado» en modos de frenado regenerativo. Esa lectura física no aparece en el informe.
\end{enumerate}

\subsection*{Frontera con el observador: $T_l$ vs.\ $T_{ld}$ \evaluar}
El modelo mecánico del observador usa $T_l(t)/r$ como perturbación residual, mientras que la planta equivalente del controlador externo, tras la compensación de gravedad del \emph{precomputed torque}, usa $T_{ld}(t)/r$. Dado que $T_l = g\,k_l\sin(\theta_l)+T_{ld}$ (guía, Ec.~2), si la gravedad ya está compensada el residuo coherente en el observador sería $T_{ld}/r$. \textbf{No se modificó} por estar en la zona delicada del ID~22 (qué señal entra al observador); conviene decidirlo junto con la unificación de los nombres de la consigna de torque.

% ===================================================================
\section{Notas técnicas}
\begin{itemize}
  \item $R_{s0}=1{,}058\ \Omega$ \textbf{no se eliminó}: es la resistencia congelada a la temperatura media de operación ($\approx 29{,}5\,^\circ$C) que define el modelo LTI (42 usos); es distinta de $R_{sREF}=1{,}02\ \Omega$ @ $20\,^\circ$C y está definida explícitamente en el texto. Solo se corrigió su uso indebido dentro de la ley física de variación térmica (ID 10).
  \item Los \emph{listings} de MATLAB no se modificaron (los nombres \texttt{lambda\_m}, \texttt{i\_ds}, etc.\ son variables del código, no símbolos del informe), salvo la nota de correspondencia del ID 23.
  \item En la presentación se replicaron todos los cambios aplicables (mismos IDs), manteniendo el espejo informe$\leftrightarrow$presentación.
  \item Tipografía menor pendiente de unificar: la sección del controlador externo usa punto decimal ($0.0396$, $31.656$, $0.75$) mientras buena parte del informe usa coma decimal ($0{,}75$).
  \item Archivos de respaldo en el directorio del proyecto: \texttt{main\_backup2.tex} y \texttt{main\_presentacion\_backup2.tex} (estado previo al lote principal), y \texttt{aux\_main\_anterior.tex} (\textbf{versión del informe subida por el compañero}, punto de partida de toda esta comparación).
\end{itemize}

\end{document}
